\begin{solution}
这类题看上去在测长度，其实真正要写的是“厚度沿长度方向怎样线性增长”。也就是说，先把劈尖的几何斜率写成 \(\alpha\approx D/L\)，再用条纹间距把 \(\alpha\) 反推出去。

这里空气劈尖取 \(n=1\)，且 \(30\) 条明纹跨越的是 \(29\) 个条纹间距，所以
\[
l=\frac{4.295}{29}\,\si{mm}\approx \SI{0.1481}{mm},\qquad
D=L\alpha=L\frac{\lambda}{2l}
=\frac{28.880\times 589.3\times10^{-6}}{2\times0.1481}\,\si{mm}
\approx \SI{0.0574}{mm}.
\]
细丝直径约为 \(\SI{5.74e-2}{mm}=\SI{57.4}{\micro m}\)，落在几十微米量级是合理的。这里不能把 \(30\) 条条纹误当成 \(30\) 个间距，也不能漏掉空气层对应 \(n=1\)。
\end{solution}

\begin{exercise}[空气劈尖条纹弯向厚端时怎样判断凸凹缺陷]
光学平板玻璃 \(A\) 与待测工件 \(B\) 间形成空气劈尖，左侧为空气膜薄端、右侧为空气膜厚端。用波长 \(\lambda=\SI{500}{nm}\) 的单色光垂直照射，反射光干涉条纹中弯曲部分向厚端偏移，且顶点恰好与其右边相邻直条纹的连线相切。判断工件上表面缺陷是凸起还是凹槽，并求最大高度或深度。
\end{exercise}
